\section{$n$阶行列式的概念}

\subsection{全排列的逆序数、对换}
把$n$个不同的元素排成一列，叫做这$n$个元素的**全排列**（简称排列）。称$n$个不同元素的排列为$n$阶排列。

为了方便，把自然数$1,2,\cdots,n$视为$n$个不同的元素的代表。用$p_i$表示这$n$个数中的一个$(i=1,2,\cdots,n)$，且当$i\neq j$时$p_i\neq p_j$，于是 $p_1 p_2 \cdots p_n$ 便是 $1, 2, \cdots, n$ 的一个排列。

对于排列 $p_1 p_2 \cdots p_n$，我们把排在 $p_i$ 前面且比 $p_i$ 大的数的个数 $t_i$ 称为全排列 $p_1 p_2 \cdots p_n$ 中 $p_i$ 的逆序数，把这个排列中各数的逆序数之和
\[
t_1 + t_2 + \cdots + t_n
\]
称为这个排列的**逆序数**，记为 $t(p_1 p_2 \cdots p_n)$。

逆序数为奇数的排列称为**奇排列**；逆序数为偶数的排列称为**偶排列**。显然，排列 $12 \cdots n$ 的逆序数为 $0$，故它是偶排列。称此排列为**自然排列**。

在一个排列中，将某两个数的位置对调（其他数不动）的变动叫做一个**对换**。两个相邻数的对换称为**相邻对换**。

\begin{theorem}
一个排列中的任意两个数对换后，排列改变奇偶性。
\end{theorem}

\subsection{$n$阶行列式的概念}
\begin{definition}
设$n^2$个数，排成$n$行$n$列的数表
\[
\begin{pmatrix}
a_{11} & a_{12} & \cdots & a_{1n}\\
a_{21} & a_{22} & \cdots & a_{2n}\\
\vdots & \vdots & & \vdots\\
a_{n1} & a_{n2} & \cdots & a_{nn}
\end{pmatrix}
\]
% 此处可继续补充n阶行列式定义表达式
\end{definition}


\section{行列式的性质}

\begin{property}
行列式与它的转置行列式相等。
\end{property}

\begin{property}
互换行列式的两行（列），行列式变号。
\end{property}

\begin{corollary}
如果行列式有两行（列）完全相同，则此行列式为零。
\end{corollary}

\begin{property}
行列式一行的公因子可以提出去，或者说以数$k$乘行列式的一行相当于用这个数乘此行列式。
\end{property}

\begin{corollary}
若行列式中有两行（列）元素成比例，则行列式等于零。
\end{corollary}

\begin{property}
如果行列式的某一行是两项和的形式，那么这个行列式就等于两个行列式的和。这两个行列式，除第$i$行以外的其他行与原来行列式完全相同。
\end{property}


\section{行列式的展开定理}

在给定的$n$阶子式中，把$a_{ij}$所在的第$i$行和第$j$列划去，余下的元素按原来的排法构成的$n-1$阶行列式叫做元素$a_{ij}$的**余子式**，记为$M_{ij}$；而$a_{ij}$的**代数余子式**记作$A_{ij}$，定义为
\[
A_{ij}=(-1)^{i+j}M_{ij}.
\]

% 引理位置，原文截断，保留占位
% \begin{lemma}
% 如果除$a_{ij}$以外……
% \end{lemma}

\begin{theorem}
行列式等于它的任一行（列）的各元素与其代数余子式的乘积之和，即
\[
D=\sum_{j=1}^{n}a_{ij}A_{ij}\quad(i=1,2,\dots,n)
\]
\[
D=\sum_{i=1}^{n}a_{ij}A_{ij}\quad(j=1,2,\dots,n)
\]
称本定理为行列式按行（列）展开法则。
\end{theorem}

\begin{theorem}
行列式$D$的任一行（列）元素与另一行（列）对应元素的代数余子式乘积之和为零。
\[
\sum_{k=1}^{n}a_{ik}A_{jk}=0\quad(i\neq j),\qquad
\sum_{k=1}^{n}a_{ki}A_{kj}=0\quad(i\neq j).
\]
\end{theorem}


\section{Cramer法则}

设有含$n$个未知数的线性方程组
\[
\begin{cases}
a_{11}x_1+a_{12}x_2+\cdots+a_{1n}x_n=b_1\\
a_{21}x_1+a_{22}x_2+\cdots+a_{2n}x_n=b_2\\
\quad\vdots\\
a_{n1}x_1+a_{n2}x_2+\cdots+a_{nn}x_n=b_n
\end{cases}
\]
% 后续Cramer法则内容自行续写